
We organize related work around two threads that our work bridges: (A) weight-space encoders for model
property prediction, and (B) generalization gap characterization.

\subsubsection{2.1 Weight-Space Encoders for Model Property Prediction}

\textbf{Flat baselines and the model zoo benchmark.} Unterthiner et al. [2020] introduced the foundational
benchmark: a zoo of ~14,000 small CNNs trained on CIFAR-10 with varying hyperparameters, with test accuracy
labels. A position-indexed MLP operating on flattened, sorted weight tensors achieves Spearman r $\approx$ 0.85 on
test accuracy prediction, establishing the competitive bar. Eilertsen et al. [2020] complemented this with
hand-crafted weight statistics (spectral norms, layer means, singular value distributions), showing that
informative properties can be extracted without learned encoders. Both works treat test accuracy as the
sole prediction target.

\textbf{Permutation-equivariant encoders.} The core insight motivating equivariant architectures is that neural
network weights are only defined up to permutation of hidden neurons: any permutation of neurons within a
layer, with corresponding column/row permutations of adjacent weight matrices, yields a functionally
equivalent network. Flat MLPs break this symmetry by sorting weights --- a preprocessing heuristic that may
introduce spurious features. Navon et al. [2023] proposed Deep Weight Space (DWS) networks: a family of
permutation-equivariant layers that compute statistics shared across neuron permutation equivalence classes.
DWS achieves Spearman r $\approx$ 0.9 on Unterthiner test accuracy prediction, substantially outperforming the flat
baseline. Crucially, DWS enforces within-layer equivariance --- each layer's neurons are treated as an
interchangeable set --- but does not explicitly model cross-layer interactions.

\textbf{Cross-layer attention.} Zhou et al. [2023] introduced Neural Functional Transformers (NFT): a
transformer-based architecture that attends across the full sequence of weight matrices while respecting
the symmetry group of the neural network. By operating on sequences of weight rows and columns across
layers, NFT captures inter-layer weight co-variation --- a qualitatively different inductive bias from
DWS's within-layer averaging. NFT achieves competitive performance on test accuracy prediction and
enables symmetry-aware fine-tuning of neural networks. Our work is the first to apply NFT to gap
prediction and to demonstrate that its cross-layer attention captures gap-specific signal.

\textbf{Graph neural network encoders.} Kofinas et al. [2024] represent neural networks as graphs --- neurons
as nodes, weights as edges --- and apply standard graph neural networks to produce equivariant weight
embeddings. This representation is theoretically appealing but imposes a rigid graph connectivity
structure that may over-constrain the representation for distributed properties like generalization gap.
In our experiments, GNN achieves the lowest gap Spearman (r = 0.3747) among all tested encoders,
below even the flat baseline.

\textbf{Universal and cross-architecture encoders.} Trabucco et al. [2024] proposed Universal Neural
Functionals (UNFN) designed to process networks of varying architectures without retraining. Schürholt
et al. [2022] developed hyper-representations via self-supervised pretraining on model zoos. Our study
is restricted to a single architecture family to ensure clean controlled comparison.

\textbf{The gap blind spot.} A striking commonality across all encoder papers above is exclusive focus on
test accuracy as the prediction target. Generalization gap --- while computable from the same model zoo
metadata --- has not been systematically studied as a prediction objective, and target-specificity
has never been examined. Our work addresses this gap directly.

\subsubsection{2.2 Generalization Gap Characterization}

\textbf{Complexity measures.} A parallel literature characterizes generalization gap using hand-crafted
complexity measures: sharpness of the loss landscape, margin-based bounds, PAC-Bayes flatness measures.
Jiang et al. [2020] conducted a comprehensive empirical study comparing 40 complexity measures as
predictors of generalization gap across model zoos, finding that no single measure dominates across
settings. Our contribution adds learned weight-space encoders to this comparison family, showing that
NFT achieves Spearman r > 0.5 --- comparable to the better complexity measures in their study, without
requiring problem-specific feature engineering.

\textbf{Theoretical connections.} PAC-Bayes generalization bounds are phrased in terms of the KL divergence
between learned and prior weight distributions --- a quantity that is computed over the full weight tensor
and is permutation-invariant by construction. This theoretical connection motivates the hypothesis that
permutation-equivariant encoders may have an inductive bias aligned with gap signal. Our experimental
findings partially support this (NFT captures genuine gap-specific signal) while also revealing that the
connection is not simply ''equivariance $\rightarrow$ gap advantage'' --- architecture specifics (cross-layer vs.
within-layer) matter.

\subsubsection{2.3 Positioning Our Work}

Our work differs from prior encoder papers in three ways: (1) we treat generalization gap as the primary
prediction target rather than a secondary consideration; (2) we conduct controlled dual-target evaluation
that directly tests target-specificity; (3) we introduce partial Spearman analysis to decompose
gap-specific vs. test-accuracy-correlated weight-space signal. We differ from complexity measure papers
by using learned encoders (no feature engineering) and by focusing on encoder architecture comparison.

